\chapter{نظرية الحركة والغاز المثالي}\label{ch:b1:kinetic-theory}

هواء قاعة الدرس يزن ما يزنه رجل بالغ، ويحوي من
الجزيئات أكثر مما على الأرض من حبات الرمل، يطير كلٌّ منها بسرعة
رصاصة بندقية ويصطدم عشرة مليارات مرة في الثانية
بجيرانه. ولا أحد يستطيع تتبع واحد منها؛ ولا أحد يحتاج إلى ذلك. فحفنة
من المتوسطات --- \omterm{def:b1:kinetic-theory:pressure}{الضغط} \omterm{def:b1:kinetic-theory:temperature}{ودرجة الحرارة} والحجم والكمية --- تصف
الهواء وصفًا تامًّا من أجل كل غرض هندسي، وصورة بسيطة
لجزيئات ترتدّ عن الجدران تفسّر من أين تأتي تلك المتوسطات
وكيف تترابط. ويفتتح هذا الفصل جزء الترموديناميك
من هذا المجلد بالنظر إلى المادة من طرفيها: الفوضى المجهرية
والهدوء العياني، ونظرية الحركة التي تربط بينهما
في حالة \omterm{def:b1:kinetic-theory:model}{الغاز المثالي}.

\section{مستويان للوصف}

\begin{definition}[المستويان المجهري والعياني]\label{def:b1:kinetic-theory:scales}
\index{مستوى مجهري}\index{مستوى عياني}\index{ثابت أفوغادرو}\index{مول}
في المستوى \emph{المجهري} تتكوّن المادة من جزيئات: ففي $\qty{1}{mol}$
منها $N_A = \qty{6.02e23}{}$ جزيئًا، قطر كلٍّ منها نحو \qty{0.3}{nm}؛
وفي غاز عند الشروط العادية تتباعد \qty{3}{nm} وتتحرك بمئات
الأمتار في الثانية وتتصادم كل \qty{70}{nm} تقريبًا.
وفي المستوى \emph{العياني} تُوصف جملة ببضعة
\emph{متغيرات حالة} --- \omterm{def:b1:kinetic-theory:pressure}{الضغط} $P$ والحجم $V$ \omterm{def:b1:kinetic-theory:temperature}{ودرجة الحرارة} $T$
وكمية المادة $n$ --- وهي متوسطات على أعداد هائلة من
الجزيئات. ويكون متغير ما \emph{امتداديًّا} إذا تضاعف حين تتضاعف
الجملة ($V$ و $n$ والكتلة والطاقة)، و\emph{شدّيًّا} إذا لم يتضاعف
($P$ و $T$ والكثافة).
\end{definition}

\begin{definition}[التوازن الترمودينامي]\label{def:b1:kinetic-theory:equilibrium}
\index{توازن ترمودينامي}\index{متغير حالة}
تكون جملة في \emph{توازن ترمودينامي} حين تكون متغيرات حالتها
منتظمة وثابتة في الزمن ولا يعبرها أيّ تدفق عياني للمادة أو
للطاقة. وعندئذ فقط يكون $P$ و $T$ معرَّفين للجملة
كلها؛ و\emph{معادلة الحالة} تربطهما بالمقدارين $V$ و $n$.
\end{definition}

\begin{definition}[الضغط]\label{def:b1:kinetic-theory:pressure}
\index{ضغط}\index{باسكال}\index{بار}
يدفع مائع ساكن كل عنصر سطحي $\dd S$ من جدار بقوة
$\dd\vect F = P\,\dd S\,\vect n$ ناظمية عليه، نحو الجدار؛ و $P$ هو
\emph{الضغط}، بالباسكال ($\qty{1}{Pa} = \qty{1}{N/m^2}$)؛ و $\qty{1}{bar}
= \qty{1e5}{Pa}$ و $\qty{1}{atm} = \qty{1.013e5}{Pa}$. وعند نقطة معطاة لا
يتوقف على توجّه السطح (\cref{ch:b1:fluid-statics}).
\end{definition}

\begin{definition}[درجة الحرارة]\label{def:b1:kinetic-theory:temperature}
\index{درجة حرارة}\index{كلفن}\index{المبدأ الصفري}
جملتان متلامستان عبر جدار يسمح بمرور الطاقة تبلغان، بعد
حين، حالة مشتركة: فتكونان عندئذ عند \emph{درجة الحرارة} نفسها؛
وجسمان كلٌّ منهما في توازن مع ثالث يكونان في توازن أحدهما
مع الآخر (وهو \emph{المبدأ الصفري}). ويُثبَّت سلّم كلفن بالقيمة
$k_B = \qty{1.380649e-23}{J/K}$ بالضبط؛ وبالسيلزيوس: $\theta = T - 273.15$.
ونظرية الحركة أدناه تعطي $T$ معناه المجهري.
\end{definition}

\section{النموذج الحركي للغاز المثالي}

\begin{definition}[الغاز المثالي، النموذج المجهري]\label{def:b1:kinetic-theory:model}
\index{غاز مثالي}\index{غاز تام}
\emph{الغاز المثالي} (أو الغاز التام) مجموعة من $N$ جزيئًا،
تُعامَل كنقاط كتلتها $m$، تتحرك بحرّية وعشوائيًّا --- بلا
تآثر إلا في تصادمات قصيرة --- في وعاء حجمه $V$.
وعند التوازن يكون توزيع السرعات متماثل المناحي ومستقرًّا؛
و $n^* = N/V$ هي الكثافة العددية و $\langle v^2\rangle$ متوسط مربع
السرعة.
\end{definition}

\begin{theorem}[الضغط الحركي]\label{thm:b1:kinetic-theory:pressure}
\index{ضغط حركي}
\omterm{def:b1:kinetic-theory:pressure}{الضغط} الذي يؤثر به الغاز على الجدران هو
\[
P = \tfrac13\,n^*m\langle v^2\rangle = \tfrac23\,n^*\langle\epsilon_k\rangle ,
\]
حيث $\langle\epsilon_k\rangle = \tfrac12 m\langle v^2\rangle$ هو متوسط
\omterm{thm:b1:work-and-energy:ke}{الطاقة الحركية} لجزيء.
\end{theorem}

\begin{proof}[برهان جزئي]
خذ النموذج المبسَّط الذي تكون فيه سرعة الجزيئات كلها $u$ وتتحرك
على امتداد المحاور الثلاثة، سدسها في كل اتجاه. فالجزيء
الذي يصطدم بالجدار $x = 0$ مباشرةً يرتدّ ارتدادًا مرنًا ويسلّمه
\omterm{def:b1:newton-dynamics:momentum}{كمية الحركة} $2mu$. وفي الزمن $\dd t$، تكون الجزيئات المتجهة نحو
الجدار التي تبلغ مساحة $S$ هي الواقعة في شريحة سماكتها $u\,\dd t$،
أي سدس الجزيئات $n^*Su\,\dd t$ التي فيها. والقوة $=$ 
\omterm{def:b1:newton-dynamics:momentum}{كمية الحركة} في وحدة الزمن $= (n^*Su\,\dd t/6)(2mu)/\dd t = \tfrac13 n^*mu^2S$، ومنه
$P = \tfrac13 n^*mu^2$. والتوسيط على التوزيع الحقيقي للسرعات
والاتجاهات يحوّل $u^2$ إلى $\langle v^2\rangle$ (إذ يعطي تماثل المناحي
$\langle v_x^2\rangle = \tfrac13\langle v^2\rangle$)؛ ويفعل ذلك مجلد السنة الثانية
بالكامل.
\end{proof}

\begin{omfigure}
\begin{tikzpicture}[font=\small, x=1cm, y=1cm]
  \draw[thick] (0,-1.6) -- (0,1.6); \fill[pattern=north east lines] (-0.2,-1.6) rectangle (0,1.6);
  \draw[dashed] (1.8,-1.6) -- (1.8,1.6);
  \draw[Stealth-Stealth] (0,-1.85) -- (1.8,-1.85) node[midway, below] {$u\,\dd t$};
  \foreach \p in {(0.6,1.1),(1.2,0.5),(0.9,-0.2),(1.5,-0.9),(0.4,-1.2),(1.3,1.3)} {\fill[omDef] \p circle (2pt);}
  \draw[-Stealth, omDef, thick] (1.2,0.5) -- (0.5,0.5);
  \draw[-Stealth, omThm, thick] (0.25,0.2) -- (0.95,0.2);
  \node[omDef, font=\footnotesize, above] at (0.85,0.55) {$m\vect u$};
  \node[omThm, font=\footnotesize, below] at (0.6,0.15) {$-m\vect u$};
  \node[right, font=\footnotesize] at (2.0,0.9) {كمية الحركة $2mu$ المسلَّمة};
  \node[right, font=\footnotesize] at (2.0,0.55) {إلى الجدار في كل ارتطام};
  \node[right, font=\footnotesize] at (2.0,-0.5) {$\tfrac16 n^*Su\,\dd t$ ارتطامًا في $\dd t$};
  \begin{scope}[xshift=7.2cm, yshift=-1.6cm]
  \begin{axis}[omaxis, axis x line=bottom, axis y line=left, width=6.4cm, height=4.2cm, xlabel={$v$ (\unit{m/s})}, ylabel={$f(v)$},
      xlabel style={at={(axis description cs:0.5,-0.1)}, anchor=north},
      xmin=0, xmax=1400, ymin=0, ymax=1.2, xtick={500,1000}, ytick=\empty, anchor=origin, at={(0,0)}]
    \addplot[omDef, thick, domain=0:1400, samples=120] {(x/422)^2*exp(1-(x/422)^2)};
    \addplot[omThm, thick, domain=0:1400, samples=120] {0.707*(x/597)^2*exp(1-(x/597)^2)};
    \node[omDef, font=\footnotesize] at (axis cs:430,1.12) {\qty{300}{K}};
    \node[omThm, font=\footnotesize] at (axis cs:900,0.72) {\qty{600}{K}};
  \end{axis}
  \end{scope}
\end{tikzpicture}

{\small يسارًا: الأصل الحركي \omterm{def:b1:kinetic-theory:pressure}{للضغط} --- فالجزيئات التي تضرب
جدارًا وترتدّ تسلّمه كمية حركة؛ وعددها في الثانية
مضروبًا في $2mu$ هو القوة. ويمينًا: السرعات الحقيقية موزَّعة (بتوزيع
ماكسويل، وهنا من أجل الآزوت عند درجتَي حرارة)؛ وتحلّ \omterm{def:b1:kinetic-theory:kinetictemp}{السرعة التربيعية المتوسطة}
محلّ $u$ الوحيدة في النموذج.}
\end{omfigure}

\begin{definition}[درجة الحرارة الحركية؛ معادلة الحالة]\label{def:b1:kinetic-theory:kinetictemp}
\index{درجة حرارة حركية}\index{معادلة الحالة!للغاز المثالي}\index{ثابت الغازات}\index{سرعة تربيعية متوسطة}
تُعرَّف \omterm{def:b1:kinetic-theory:temperature}{درجة حرارة} \omterm{def:b1:kinetic-theory:model}{غاز مثالي} \emph{تعريفًا} بمتوسط \omterm{thm:b1:work-and-energy:ke}{الطاقة
الحركية} لجزيئاته،
\[
\langle\epsilon_k\rangle = \tfrac12 m\langle v^2\rangle = \tfrac32 k_BT ,
\]
بحيث، بالجمع مع \cref{thm:b1:kinetic-theory:pressure}،
\[
PV = Nk_BT = nRT, \qquad R = N_Ak_B = \qty{8.314}{J/(mol.K)} ,
\]
وهي \emph{معادلة حالة \omterm{def:b1:kinetic-theory:model}{الغاز المثالي}}. و\emph{السرعة التربيعية
المتوسطة} هي $u = \sqrt{\langle v^2\rangle} = \sqrt{3k_BT/m} = \sqrt{3RT/M}$،
حيث $M$ الكتلة المولية.
\end{definition}

\begin{example}[أعداد من أجل الهواء]\label{ex:b1:kinetic-theory:air}
الآزوت عند \qty{300}{K}: $u = \sqrt{3 \times 8.314 \times 300/0.028} =
\qty{517}{m/s}$؛ والهدروجين، $\qty{1930}{m/s}$؛ وكلما ثقُل الغاز بطُؤ، وفق
$1/\sqrt M$. وعند \qty{1}{bar} و \qty{300}{K} يكون $n^* = P/k_BT = \qty{2.4e25}{m^{-3}}$:
أي \qty{24}{L} لكل \omterm{def:b1:kinetic-theory:scales}{مول}، و \qty{2.7e22}{} جزيئًا في اللتر. وقاعة درس حجمها
\qty{200}{m^3}: أي \qty{8300}{mol} و \qty{240}{kg} من الهواء.
\end{example}

\begin{remark}[لماذا يعمل النموذج، ومتى يفشل]\label{rem:b1:kinetic-theory:validity}
\index{متوسط المسار الحرّ}
بين تصادمين يطير جزيء هواء \emph{متوسط مسار حرّ}
$\lambda \approx 1/(\sqrt2\pi d^2n^*) \approx \qty{70}{nm}$ --- أي مئتَي
ضعف حجمه: فالغاز خال في معظمه، والتآثرات نادرة،
ويصحّ $PV = nRT$ بأفضل من واحد في المئة عند الضغوط العادية. وعند \omterm{def:b1:kinetic-theory:pressure}{الضغط}
العالي أو \omterm{def:b1:kinetic-theory:temperature}{درجة الحرارة} المنخفضة يصير حجم الجزيئات نفسه وتجاذبها مهمَّين:
فمعادلة فان دير فالس $(P + an^2/V^2)(V - nb) = nRT$
تصحّح للأمرين، ودون \omterm{def:b1:kinetic-theory:temperature}{درجة حرارة} حرجة يتسيّل الغاز
(\cref{ch:b1:phase-changes}).
\end{remark}

\begin{omfigure}
\begin{tikzpicture}
\begin{axis}[omaxis, axis x line=bottom, axis y line=left, width=10cm, height=6cm, xlabel={$V$}, ylabel={$P$},
    xlabel style={at={(axis description cs:0.5,-0.08)}, anchor=north},
    xmin=0, xmax=5.2, ymin=-0.4, ymax=4.2, xtick=\empty, ytick=\empty]
  \addplot[omDef, thick, domain=0.3:5.2, samples=100] {1.0/x};
  \addplot[omDef, thick, domain=0.45:5.2, samples=100] {1.5/x};
  \addplot[omDef, thick, domain=0.6:5.2, samples=100] {2.0/x};
  \addplot[omThm, thick, domain=0.75:5.2, samples=100] {2.5/x};
  \node[omDef, font=\footnotesize, below] at (axis cs:4.9,0.19) {$T_1$};
  \node[omThm, font=\footnotesize, above] at (axis cs:4.9,0.53) {$T_4$};
  \draw[-Stealth, omExo] (axis cs:1.3,1.0) -- (axis cs:2.3,2.0) node[above right, font=\footnotesize] {$T$ متزايدة};
  \node[font=\footnotesize] at (axis cs:3.4,3.2) {$PV = nRT$: متساويات الحرارة $T_1 < \dots < T_4$};
\end{axis}
\end{tikzpicture}

{\small مخطط كلابيرون \omterm{def:b1:kinetic-theory:model}{لغاز مثالي}: عند كل \omterm{def:b1:kinetic-theory:temperature}{درجة حرارة} يكون
خط تساوي الحرارة قطعًا زائدًا $P = nRT/V$، أعلى كلما ارتفع $T$. والتحول
مسار في هذا المستوي؛ والدورة عروة مغلقة (\cref{ch:b1:heat-engines}).}
\end{omfigure}

\section{الطاقة الداخلية للغاز المثالي}

\begin{definition}[الطاقة الداخلية]\label{def:b1:kinetic-theory:U}
\index{طاقة داخلية}
\emph{الطاقة الداخلية} $U$ لجملة هي مجموع الطاقات الحركية
لجزيئاتها في المرجع الذي تسكن فيه الجملة، ومجموع
طاقاتها الكامنة المتبادلة؛ وهي دالة حالة امتدادية.
\end{definition}

\begin{proposition}[الطاقة الداخلية للغازات المثالية]\label{prop:b1:kinetic-theory:Ugas}
\index{غاز أحادي الذرة}\index{غاز ثنائي الذرة}\index{سعة حرارية بحجم ثابت}\index{قانونا جول}
من أجل \omterm{def:b1:kinetic-theory:model}{غاز مثالي} تتوقف $U$ على $T$ وحدها (وهو قانون جول الأول):
\begin{gather*}
U = \tfrac32 nRT \quad (\text{أحادي الذرة: He, Ne, Ar}), \\
U = \tfrac52 nRT \quad (\text{ثنائي الذرة عند درجات الحرارة العادية: N}_2, \text{O}_2, \text{H}_2),
\end{gather*}
ومن ثمّ تكون السعة الحرارية المولية بحجم ثابت $C_{V,m} = \dd U_m/\dd T$
هي $\tfrac32 R = \qty{12.5}{J/(mol.K)}$ و $\tfrac52 R = \qty{20.8}{J/(mol.K)}$
على الترتيب.
\end{proposition}

\begin{proof}[برهان جزئي]
لا \omterm{def:b1:work-and-energy:conservative}{طاقة كامنة} متبادلة بالفرض، ومنه $U = N\langle\epsilon_k\rangle
= \tfrac32 Nk_BT$ من أجل جزيئات نقطية. والجزيء الثنائي الذرة يدور أيضًا:
فتحمل كلٌّ من درجتَي حريته الدورانيتين، عند التوازن،
المقدار $\tfrac12 k_BT$ نفسه الذي تحمله كل درجة انسحابية (وهي مبرهنة
\emph{التوزّع المتساوي}، مقبولةً؛ ويشتقّها مجلد السنة الثالثة ويفسّر لماذا
لا ينضمّ الاهتزاز إلا فوق ألف كلفن).
\end{proof}

\begin{example}[ما يعادله هواء غرفة]\label{ex:b1:kinetic-theory:room}
من أجل \qty{8300}{mol} من الهواء في قاعة الدرس عند \qty{300}{K}: $U = \tfrac52 \times
8300 \times 8.314 \times 300 = \qty{5.2e7}{J}$ --- أي طاقة لتر ونصف
من البنزين، محفوظة في حركة جزيئية. وتسخين القاعة بمقدار $1$ \omterm{def:b1:kinetic-theory:temperature}{كلفن}
يكلّف $\tfrac52 nR = \qty{170}{kJ}$ (بحجم ثابت).
\end{example}

\section{الأطوار المتكاثفة}

\begin{proposition}[نموذج الطور غير القابل للانضغاط ولا للتمدد]\label{prop:b1:kinetic-theory:condensed}
\index{طور متكاثف}\index{تمدد حراري}\index{انضغاطية}
تُنمذَج السوائل والأجسام الصلبة، في هذا المستوى، أطوارًا ثابتة
الحجم --- لا تنضغط ولا تتمدد --- تتوقف طاقتها الداخلية
على $T$ وحدها: $\dd U = C\,\dd T$ مع $C = mc$، حيث $c$ السعة الحرارية
الكتلية (\qty{4180}{J/(kg.K)} للماء و \qty{450}{J/(kg.K)} للفولاذ). أما
الانحرافات الحقيقية الصغيرة فتُقاس بمعامل \omterm{prop:b1:kinetic-theory:condensed}{التمدد الحراري}
$\alpha = (1/V)(\partial V/\partial T)_P$ ($\sim\qty{1e-5}{K^{-1}}$
للمعادن، و \qty{2e-4}{K^{-1}} للسوائل) وبالانضغاطية
المتساوية الحرارة $\chi_T = -(1/V)(\partial V/\partial P)_T$ ($\sim\qty{5e-10}{Pa^{-1}}$
للماء، مقابل $1/P \approx \qty{1e-5}{Pa^{-1}}$ لغاز).
\end{proposition}
\admitted

\begin{example}[كم هو عصيّ على الانضغاط]\label{ex:b1:kinetic-theory:water}
الماء في قاع خندق ماريانا (\qty{1100}{bar}) ينضغط
بالمقدار $\chi_T\Delta P \approx 5\times10^{-10} \times 1.1\times10^8 = 5\%$؛ وسكة
فولاذية طولها \qty{30}{m} تُسخَّن بمقدار $50$ \omterm{def:b1:kinetic-theory:temperature}{كلفن} تطول بالمقدار $\alpha L\Delta T =
\qty{18}{mm}$ --- ومن هنا فواصل التمدد. وأمام غاز ينتصف
حجمه تحت \omterm{def:b1:kinetic-theory:pressure}{ضغط} مضاعف، يكون نموذج الحجم الثابت
جيدًا إلى بضعة في المئة في معظم أوضاع هذا المجلد.
\end{example}

\section{تمارين}

\begin{exercise}[$\star$]\label{exo:b1:kinetic-theory:1}
احسب عدد الجزيئات في \qty{1.0}{L} من غاز عند \qty{0}{\degreeCelsius} و
\qty{1}{atm}؛ والمسافة المتوسطة بين الجيران؛ وقارنها بحجم الجزيء
\qty{0.3}{nm}.
\end{exercise}

\begin{exercise}[$\star$]\label{exo:b1:kinetic-theory:2}
احسب السرعات التربيعية المتوسطة للغازات H$_2$ و N$_2$ و O$_2$ و CO$_2$ عند \qty{300}{K}،
ثم من أجل N$_2$ عند \qty{77}{K} (الآزوت المغلي).
\end{exercise}

\begin{exercise}[$\star$]\label{exo:b1:kinetic-theory:3}
يُنفخ إطار سيارة حجمه \qty{30}{L} إلى \omterm{def:b1:kinetic-theory:pressure}{ضغط} مقياسي مقداره
\qty{2.5}{bar} عند \qty{20}{\degreeCelsius}. احسب كمية الهواء وكتلته في الداخل؛
\omterm{def:b1:kinetic-theory:pressure}{والضغط} المطلق بعد قيادة تسخّنه إلى \qty{50}{\degreeCelsius}.
\end{exercise}

\begin{exercise}[$\star$]\label{exo:b1:kinetic-theory:4}
قدّر \omterm{rem:b1:kinetic-theory:validity}{متوسط المسار الحرّ} لجزيئات الهواء عند \qty{1}{bar} و \qty{300}{K}
(مع $d = \qty{0.37}{nm}$)، وتواتر تصادم جزيء واحد. وعند
أيّ \omterm{def:b1:kinetic-theory:pressure}{ضغط} يبلغ \omterm{rem:b1:kinetic-theory:validity}{متوسط المسار الحرّ} \qty{1}{m} (أي ``خلاء'')؟
\end{exercise}

\begin{exercise}[$\star\star$]\label{exo:b1:kinetic-theory:5}
احسب كثافة الهواء عند \qty{1.013}{bar} و \qty{20}{\degreeCelsius}
(مع $M = \qty{29}{g/mol}$)، ثم عند \qty{100}{\degreeCelsius}. واحسب حمولة
منطاد هواء ساخن حجمه \qty{2800}{m^3} (أي فرق وزنَي الهواء).
\end{exercise}

\begin{exercise}[$\star\star$]\label{exo:b1:kinetic-theory:6}
الهواء $21\%$ منه O$_2$ بالجزيئات. احسب \omterm{def:b1:kinetic-theory:pressure}{الضغط} الجزئي للأكسجين عند سطح البحر؛
وعند \qty{5000}{m} حيث $P = \qty{0.54}{atm}$؛ وعند قمة إفرست
(\qty{0.33}{atm}). ولماذا يعاني الجسم هناك؟
\end{exercise}

\begin{exercise}[$\star\star$]\label{exo:b1:kinetic-theory:7}
\omterm{def:b1:kinetic-theory:scales}{مول} واحد من CO$_2$ عند \qty{300}{K} في \qty{1.0}{L}: احسب \omterm{def:b1:kinetic-theory:pressure}{الضغط} بقانون
\omterm{def:b1:kinetic-theory:model}{الغاز المثالي}، ثم بفان دير فالس ($a = \qty{0.364}{Pa.m^6/mol^2}$ و
$b = \qty{4.27e-5}{m^3/mol}$). فأيّ التصحيحين يغلب هنا؟
\end{exercise}

\begin{exercise}[$\star\star$]\label{exo:b1:kinetic-theory:8}
احسب \omterm{def:b1:kinetic-theory:U}{الطاقة الداخلية} لهواء غرفة حجمها \qty{50}{m^3} عند \qty{1}{bar} و
\qty{300}{K}؛ والطاقة اللازمة لتسخينه بمقدار \qty{5}{K} بحجم ثابت؛ والزمن
بمدفأة \qty{2}{kW} (بلا ضياعات).
\end{exercise}

\begin{exercise}[$\star\star$]\label{exo:b1:kinetic-theory:9}
سكة فولاذية طولها \qty{30}{m} ($\alpha = \qty{1.2e-5}{K^{-1}}$) بين
\qty{-10}{\degreeCelsius} و \qty{40}{\degreeCelsius}: احسب تغيّر الطول.
وحلقة نحاسية ($\alpha = \qty{1.9e-5}{K^{-1}}$) قطرها الداخلي
\qty{49.95}{mm} يجب أن تنزلق فوق محور قطره \qty{50.00}{mm}: فبكم يجب
تسخينها؟
\end{exercise}

\begin{exercise}[$\star\star\star$]\label{exo:b1:kinetic-theory:10}
احسب الطاقة لكل جزيء عند \qty{300}{K} (أي $\tfrac32 k_BT$) بالجول وبوحدة \unit{eV}؛
والسعتين الحراريتين الموليتين $C_{V,m}$ للأرغون وللآزوت؛ ولماذا تكون سعة
الآزوت أكبر مع أن جزيئاته أخفّ؟
\end{exercise}

\begin{exercise}[$\star\star\star$]\label{exo:b1:kinetic-theory:11}
أعِد اشتقاق \omterm{thm:b1:kinetic-theory:pressure}{الضغط الحركي} في نموذج الاتجاهات الستة،
ثم بيّن أنه من أجل توزيع متماثل المناحي للسرعات يكون $\langle v_x^2\rangle
= \tfrac13\langle v^2\rangle$ وأن صيغة \omterm{def:b1:kinetic-theory:pressure}{الضغط} لا تتغير.
\end{exercise}

\begin{exercise}[$\star\star\star$]\label{exo:b1:kinetic-theory:12}
قارن السرعتين التربيعيتين المتوسطتين للغازين H$_2$ و O$_2$ عند \qty{300}{K}
\omterm{prop:b1:central-forces:escape}{بسرعة الإفلات} من الأرض، وبسرعة الإفلات من القمر (\qty{2.4}{km/s}).
ومع العلم بأن نسبة الجزيئات الأسرع من $3u$ صغيرة لكنها ليست
معدومة، فسّر لماذا فقدت الأرض هدروجينها وفقد القمر غلافه
الجوي كله.
\end{exercise}

\section{مسألة: هواء القاعة}

\begin{problem}[{مسألة نهاية الأسبوع --- قاعة درس مملوءة هواءً: عدّ
جزيئاته، ووزنه، وتوقيت طيرانها، واستعادة حمولة منطاد
ورقّة هواء الجبال وسرعة الصوت من النموذج
نفسه}]
\label{pb:b1:kinetic-theory:1}
قاعة درس أبعادها $10 \times 8 \times \qty{2.5}{m}$؛ وهواؤها عند
$P = \qty{1.013e5}{Pa}$ و $T = \qty{293}{K}$، وكتلته المولية $M = \qty{29.0}{g/mol}$،
و $79\%$ منه N$_2$ و $21\%$ منه O$_2$ بالجزيئات. ولدينا $R = \qty{8.314}{J/(mol.K)}$
و $k_B = \qty{1.38e-23}{J/K}$ و $N_A = \num{6.02e23}$ و $g = \qty{9.81}{m/s^2}$.

\medskip
\noindent\textbf{الجزء I --- العدّ والوزن.}
\begin{enumerate}
  \item احسب كمية الهواء (بالمولات) وعدد الجزيئات في القاعة.
  \item احسب كتلة الهواء؛ وقارنها بشخص.
  \item احسب الكثافة العددية $n^*$ والمسافة المتوسطة بين الجزيئات
        ($n^{*-1/3}$)؛ وقارنها بقطر الجزيء \qty{0.37}{nm}.
  \item احسب \omterm{def:b1:kinetic-theory:kinetictemp}{السرعة التربيعية المتوسطة} للغاز N$_2$ وللغاز O$_2$؛ ولماذا يتحرك الغاز
        الأثقل أبطأ عند \omterm{def:b1:kinetic-theory:temperature}{درجة الحرارة} نفسها؟
  \item احسب متوسط \omterm{thm:b1:work-and-energy:ke}{الطاقة الحركية} لجزيء واحد؛ \omterm{thm:b1:work-and-energy:ke}{والطاقة الحركية}
        الانسحابية الكلية لهواء القاعة؛ وقارنها \omterm{thm:b1:work-and-energy:ke}{بالطاقة الحركية}
        لسيارة عند \qty{100}{km/h}.
  \item احسب \omterm{rem:b1:kinetic-theory:validity}{متوسط المسار الحرّ} $\lambda \approx 1/(\sqrt2\pi d^2n^*)$ ومتوسط
        الزمن بين تصادمَي جزيء.
\end{enumerate}

\medskip
\noindent\textbf{الجزء II --- \omterm{def:b1:kinetic-theory:pressure}{الضغط} من الارتطامات.}
\begin{enumerate}[resume]
  \item باستعمال نموذج الاتجاهات الستة مع $u = \qty{500}{m/s}$، احسب
        عدد الارتطامات في الثانية على \qty{1}{cm^2} من الجدار.
  \item احسب \omterm{def:b1:newton-dynamics:momentum}{كمية الحركة} المسلَّمة في الثانية، وتحقق من أنها
        تعيد إنتاج \omterm{def:b1:kinetic-theory:pressure}{الضغط} الجوي.
  \item احسب القوة الكلية للهواء على جدار أبعاده \qty{8}{m} في \qty{2.5}{m}؛
        ولماذا لا يتحرك الجدار؟
  \item تُسخَّن القاعة إلى \qty{303}{K} ونوافذها مغلقة
        ومحكمة: فما \omterm{def:b1:kinetic-theory:pressure}{الضغط} الجديد، وما القوة المحصّلة على الجدار الآن؟
  \item التسخين نفسه مع نافذة مفتوحة: فماذا يخرج، وكم مقداره؟
\end{enumerate}

\medskip
\noindent\textbf{الجزء III --- الدافعة.}
\begin{enumerate}[resume]
  \item احسب كثافة هواء القاعة، وكثافة هواء عند \qty{373}{K} \omterm{def:b1:kinetic-theory:pressure}{بالضغط}
        نفسه.
  \item يُملأ منطاد هواء ساخن حجمه \qty{2800}{m^3} بهواء عند
        \qty{373}{K}: احسب وزن الهواء الساخن، ووزن الهواء البارد
        المزاح، والحمولة المتاحة (بأرخميدس،
        \cref{ch:b1:fluid-statics}).
  \item يزن الغلاف والسلة والموقد والغاز \qty{300}{kg}: فكم راكبًا
        كتلة كلٍّ منهم \qty{75}{kg} يستطيع رفعهم؟
  \item إلى أيّ \omterm{def:b1:kinetic-theory:temperature}{درجة حرارة} ينبغي تسخين الهواء لمضاعفة الحمولة،
        ولماذا لا يُفعل ذلك (فالغلاف من النايلون)؟
  \item منطاد هليوم بالحجم نفسه عند \qty{293}{K}
        (مع $M = \qty{4}{g/mol}$): فما الحمولة؟ ولماذا تستعمل مناطيد
        السياحة الهواء الساخن مع ذلك؟
\end{enumerate}

\medskip
\noindent\textbf{الجزء IV --- الهواء الرقيق وسرعة الصوت.}
\begin{enumerate}[resume]
  \item احسب \omterm{def:b1:kinetic-theory:pressure}{الضغط} الجزئي للأكسجين في القاعة.
  \item عند قمة إفرست $P = \qty{0.33}{atm}$ و $T \approx \qty{250}{K}$:
        احسب الكثافة العددية لجزيئات الأكسجين وقارنها بكثافة القاعة.
  \item يهبط \omterm{def:b1:kinetic-theory:pressure}{ضغط} الغلاف الجوي مع الارتفاع تقريبًا وفق
        $P(z) = P_0\eu^{-z/H}$ مع $H = RT/Mg$ (\cref{ch:b1:fluid-statics}):
        احسب $H$ من أجل \qty{260}{K} والارتفاع الذي يكون عنده $P = P_0/2$.
  \item سرعة الصوت في \omterm{def:b1:kinetic-theory:model}{غاز مثالي} هي $c_s = \sqrt{\gamma RT/M}$
        مع $\gamma = 1.4$ للهواء (\cref{ch:b1:first-law}): احسبها عند
        \qty{293}{K} وقارنها \omterm{def:b1:kinetic-theory:kinetictemp}{بالسرعة التربيعية المتوسطة}؛ وفسّر
        قربهما.
  \item كيف تتغير $c_s$ مع $T$؟ احسب سرعة الصوت في القاعة عند
        \qty{303}{K}، وفي الهليوم عند \qty{293}{K} (مع $\gamma = 5/3$): ولماذا
        يبدو الصوت في الهليوم حادًّا؟
  \item عند قمة إفرست، هل الصوت أسرع أم أبطأ منه في
        القاعة؟ وبكم؟
  \item يُستبدل بهواء القاعة أرغون عند $P$ و $T$ نفسيهما:
        فما الذي يتغير في العدد والكتلة والسرعات \omterm{def:b1:kinetic-theory:U}{والطاقة
        الداخلية}؟
  \item احسب \omterm{def:b1:kinetic-theory:U}{الطاقة الداخلية} لهواء القاعة ($U = \tfrac52 nRT$)
        والطاقة اللازمة لتسخينه بمقدار \qty{10}{K} بحجم ثابت.
  \item لخّص في ثلاثة أسطر ما قدّمه النموذج الحركي: أيّ
        المقادير العيانية فسّرها، وانطلاقًا من أيّ متوسط مجهري
        واحد.
\end{enumerate}
\end{problem}
